\section{関数と接線}\label{節:導入変化}
これまでは横軸に時刻を取ったが,同じ計算は時間以外にも使える.\,$y=f(x)$において,\,$x$が$a$から$a+h$まで変われば,\,$y$の変化は$f(a+h)-f(a)$である.したがって
\[
 \frac{f(a+h)-f(a)}h
\]
は,横の変化に対する縦の変化の割合である.これを\textbf{平均変化率}という.二点$(a,f(a)),(a+h,f(a+h))$を結ぶ直線の傾きとしても読める.

$f(x)=x^2$,\,$a=1$としてみると,
\[
 \frac{(1+h)^2-1}h=2+h
\]
となる.これは位置から速度を求めた計算と同じである.\,$h$を$0$に近づけると傾きは$2$に近づくので,点$(1,1)$における接線は
\[
 y-1=2(x-1),\qquad y=2x-1
\]
となりそうだ.ここでは微分の厳密な定義より先に,接線の候補を求めている.

別の求め方もある.点$(1,1)$を通る傾き$m$の直線を$y=m(x-1)+1$とすると,放物線との交点は
\[
 x^2-mx+m-1=(x-1)\{x-(m-1)\}=0
\]
で決まる.二つの交点の$x$座標は$1,m-1$だから,それらが重なるのは$m=2$のときである.この場合には,二次方程式の重解という見方でも同じ直線が得られる.ただし,高次の曲線では接線が別の場所で再び交わることもある.「交点が一つだけ」という説明をそのまま一般の接線の定義にはできない.

\begin{figure}[htbp]
\centering
\begin{tikzpicture}[x=2.6cm,y=1.5cm,>=Stealth]
 \draw[->] (-0.12,0)--(2.1,0) node[right]{$x$};
 \draw[->] (0,-0.25)--(0,3.75) node[above]{$y$};
 \draw[very thick,domain=0:1.9,samples=80] plot (\x,{\x*\x}) node[right]{$y=x^2$};
 \draw[blue,thick,domain=0.45:1.75] plot (\x,{2*\x-1});
 \draw[red,thick,domain=0.55:1.8] plot (\x,{2.6*\x-1.6});
 \fill (1,1) circle(1.6pt) node[below right]{$P(1,1)$};
 \fill (1.6,2.56) circle(1.6pt) node[right]{$Q$};
 \draw[dashed] (1,0) node[below]{$1$}--(1,1);
 \draw[dashed] (1.6,0) node[below]{$1+h$}--(1.6,2.56);
 \node[blue] at (1.82,2.08) {接線};
 \node[red] at (0.65,1.6) {割線};
\end{tikzpicture}
\caption{二点を結ぶ割線の傾きを,点 $Q$ を点 $P$ に近づけて調べる.}
\label{図:導入接線}
\end{figure}
図\ref{図:導入接線}では,二点を結ぶ直線を割線と呼ぶ.点$Q$を点$P$へ近づけ,その傾きの極限を考える.

